
Our initial idea was to exploit the strict repeatability of the coins' appearance to implement a template matching algorithm.
In fact, all euro coins have a circular shape and a very specific design on both sides, which makes them easily distinguishable. 
Due to the limitations of the template matching algorithm, we decided to use all the information we had about the coins to create a cascade of classifiers that would progressively filter out non-coin objects and classify the coins.
The main features we could exploit were:
\begin{itemize}
    \item circular in shape;
    \item fixed size (diameter) for each denomination due to the fact that by hypotesis all pictures are taken from a fixed distance and resolution;
    \item color information (gold, silver, bi-metallic);
    \item texture information.
\end{itemize}
After applying some preprocessing to the input image, we could use a cascade of classifiers to detect and classify the coins.
The cascade consists of three stages:
\begin{itemize}
    \item circles detection;
    \item filtering by color;
    \item template matching.
\end{itemize}